\section{Глава 1. Предметная область и существующие решения}


\subsection{1.1 Предметная область и актуальность}
Совместное планирование поездки представляет собой последовательность взаимосвязанных действий: выбор направления и дат, формирование списка активностей, обсуждение альтернатив, распределение расходов и согласование итогового маршрута. На каждом этапе участники принимают коллективные решения, которые влияют на дальнейшие шаги. Изменение одного параметра поездки, например даты или состава участников, затрагивает сразу несколько подсистем планирования.

Ключевая проблема предметной области заключается в разрыве контекста между инструментами, когда обсуждения, маршрут и расходы фиксируются в разных сервисах. В результате участники получают несколько версий одного и того же плана, а владелец поездки вынужден вручную синхронизировать информацию. Это увеличивает вероятность ошибок, усложняет контроль договоренностей и снижает прозрачность процесса [8].

Таким образом, актуальной является задача создания единого мобильного пространства, в котором данные поездки связаны на уровне пользовательского сценария, а переход между этапами планирования выполняется без потери контекста [7][8][9].

\subsection{1.2 Описание существующих решений}
Для анализа использованы решения, представленные в проектном предложении, а также их публичные описания [8]. Рассматриваемые сервисы закрывают отдельные аспекты группового планирования и позволяют определить функциональные границы существующих подходов.

\subsubsection{1.2.1 Splitwise}
Splitwise ориентирован на учет совместных расходов и расчет взаиморасчетов между участниками группы [1]. Сервис предоставляет удобные механизмы фиксации трат и контроля задолженности, однако не покрывает задачи обсуждения идей и формирования маршрутного плана по дням в рамках единого процесса [1][8].

\subsubsection{1.2.2 Wanderlog}
Wanderlog предназначен для планирования поездок и совместного редактирования маршрута [2]. Сильной стороной решения является маршрутная часть, однако в контексте целевой постановки разрабатываемого приложения финансовые и дискуссионные сценарии представлены ограниченно и не формируют целостный контур принятия решений [2][8].

\subsubsection{1.2.3 TravelSpend}
TravelSpend специализируется на фиксации расходов во время поездки [3]. Продукт эффективен для финансового учета, но не позиционируется как единая среда, объединяющая идеи, обсуждения, маршрут и совместное управление доступом [3][8].

\subsubsection{1.2.4 TripIt}
TripIt решает задачу агрегации информации о поездке и автоматизации отдельных организационных операций [4]. При этом сервис не ориентирован на полноценную коллективную модерацию идей и их управляемый перенос в маршрут в рамках одного пользовательского пространства [4][8].

\subsubsection{1.2.5 Яндекс Путешествия}
В публикации от 3 апреля 2025 года для сервиса «Яндекс Путешествия» описан режим совместной поездки по ссылке-приглашению [5]. Данный функционал улучшает совместный доступ к информации, однако основной акцент сервиса остается на сценариях подбора и бронирования, а не на полном цикле группового планирования с интегрированным учетом обсуждений и расходов [5][8].

\subsubsection{1.2.6 RUSSPASS}
RUSSPASS ориентирован на подбор маршрутов и туристический контент [6]. По открытым описаниям сервис поддерживает информационный и навигационный контур путешествия, но не заявляет полный набор механизмов коллективного планирования, включая управляемые обсуждения и баланс расходов в рамках одной рабочей области [6][8].

Проведенный обзор показывает, что существующие продукты преимущественно специализируются на отдельных задачах. Наиболее частый подход - закрытие одного функционального сегмента (маршрут, расходы или контент), а не сквозного сценария совместной подготовки поездки.

\subsection{1.3 Описание разрабатываемого решения}
Разрабатываемое решение направлено на объединение ключевых пользовательских действий в рамках одной поездки как центральной сущности. Функциональная модель сервиса включает:
1. Авторизацию и персональный профиль пользователя.
2. Создание и управление поездками с параметрами дат, валюты и состава участников.
3. Механизм приглашений для совместного доступа.
4. Подсистему идей и обсуждений с ролями владельца и участника.
5. Маршрут по дням с операциями добавления, переноса и переупорядочивания активностей.
6. Учет расходов с балансом и статусами оплаты.
7. Погодный модуль по контексту поездки.
8. AI-подсказки с сохранением предложения в список идей.
9. Настройки уведомлений и поддержку офлайн-работы с последующей синхронизацией [7][9].

С точки зрения пользовательского результата разрабатываемое приложение формирует единую рабочую область, в которой обсуждение и фиксация решений не разрываются между внешними сервисами. Это позволяет уменьшить количество повторных действий и повысить прозрачность группового планирования [7][8].

\subsection{1.4 Сравнительный анализ существующих решений}
Для сопоставления решений использована матрица критериев, непосредственно связанных с целевой задачей разрабатываемого приложения.

\begin{table}[H]
\centering
\small
\caption{Сравнение существующих решений и разрабатываемого приложения}
\begin{tabular}{|p{2.2cm}|p{2.8cm}|p{2cm}|p{1.9cm}|p{2cm}|p{1.2cm}|p{1.4cm}|p{2.3cm}|}
\hline
Сервис & Совместное планирование в одном пространстве & Идеи и обсуждения & Маршрут по дням & Расходы и баланс & Погода & AI-подсказки & Приглашения/совместный доступ \\ 
\hline
Splitwise & Нет & Нет & Нет & Да & Нет & Нет & Да \\ 
\hline
Wanderlog & Частично & Частично & Да & Частично & Частично & Нет & Да \\ 
\hline
TravelSpend & Нет & Нет & Нет & Да & Нет & Нет & Частично \\ 
\hline
TripIt & Частично & Нет & Да & Нет & Частично & Нет & Частично \\ 
\hline
Яндекс Путешествия & Частично & Нет & Частично & Нет & Частично & Частично & Да \\ 
\hline
RUSSPASS & Частично & Нет & Частично & Нет & Частично & Нет & Частично \\ 
\hline
Разрабатываемое приложение & Да & Да & Да & Да & Да & Да & Да \\ 
\hline
\end{tabular}
\end{table}

Результаты сравнительного анализа показывают, что у рассмотренных решений присутствуют сильные стороны в отдельных функциональных доменах, однако в рамках выбранных критериев и источников [1][2][3][4][5][6] не прослеживается целостное покрытие сквозного сценария «идея -> обсуждение -> маршрут -> расходы -> синхронизация» в одном пользовательском пространстве [8].

\subsection{1.5 Функциональный разрыв и обоснование ниши разрабатываемого приложения}
По итогам анализа можно выделить функциональный разрыв, релевантный предметной области совместного планирования поездок:
1. Отсутствие единого пространства, где одновременно поддерживаются обсуждение идей, маршрутное планирование и финансовая прозрачность.
2. Недостаточная связность между этапами принятия решения и фактическим включением активности в маршрут.
3. Ограниченная поддержка коллективной работы при изменениях данных и нестабильном подключении.

Позиционирование разрабатываемого приложения определяется как интегрированный мобильный сервис для небольших групп пользователей, которым необходима единая среда принятия и фиксации решений по поездке. В отличие от специализированных решений, приложение ориентировано на целостный процесс, где каждая подсистема работает в контексте одной поездки и общей ролевой модели [7][8][9].

\subsection{1.6 Выводы по главе}
В первой главе рассмотрена предметная область совместного планирования поездок и обоснована актуальность ее цифровой интеграции. Проведен обзор существующих решений, показавший, что распространенные сервисы закрывают отдельные подзадачи, но в рамках принятых критериев сравнения не демонстрируют целостного сквозного сценария коллективной подготовки поездки.

На основе сравнительного анализа сформулирован функциональный разрыв и определено позиционирование разрабатываемого приложения как решения, объединяющего в едином пользовательском контуре управление поездкой, обсуждение идей, маршрут, расходы и вспомогательные сервисы. Полученные результаты задают основу для второй главы, в которой рассматриваются проектные решения и архитектура системы.
